% Notes of the Montreal seminar, 1970 (cote 129, typed part): archive cote 129, site page 68 (PDF page 69).
% typescript. Draft compared: none.
% AI-transcribed by Claude Opus 5.5 (claude-opus-5-5, adaptive thinking, effort high), Message Batches API, batch mtl129_full (msgbatch_01AGnDKHwy6ATCM7ehUa6263), 2026-11-13; not yet checked by eye.
\scanpage{129/68}
\mspage{\hand{III - 13}}

rentes de $p$, on montre que $A(1)$ est étale sur $S$, et on dira alors
que $A(\infty)$ est ind-étale.

Notons par ailleurs qu'il existe une notion de dualité pour
les variétés abéliennes. Si $A^*$ désigne la variété duale de la variété
abélienne $A$, on montre aisément que son groupe de BT associé $A^*(\infty)$
s'identifie au dual, en tant que groupe de BT de $A(\infty)$.

6.2. Par définition, on dit qu'un groupe de BT sur $S$ est \uline{un groupe}
\uline{de Barsotti-Tate ind-étale} si $G(1)$ est étale. Il est équivalent de dire
(les $G(n)$ admettant une suite de composition en groupes isomorphes à
$G(1)$) que tous les $G(n)$ sont étales.

Le foncteur $G \mapsto T_p(G)$ (5.4) définit une équivalence de caté-%
gories entre la catégorie des groupes de BT ind-étales et la catégorie
des \uline{faisceaux p-adiques constants tordus sans torsion} (i.e. systèmes
p-adiques de groupes finis localement libres $G_n$, étales sur $S$ et plats
sur $\Lambda_n$).

Si $S$ est connexe, la catégorie des faisceaux p-adiques cons-%
tants tordus sans torsion est bien connue: si l'on choisit un point géo-%
métrique $\overline{s} \in S$, on définit une équivalence de cette catégorie avec la
catégorie des $Z_p$-modules libres de type fini sur lesquels le groupe fon-%
damental $\pi_1(S,\overline{s})$ opère continuement en associant au système $(G_n)_{n\geq 1}$
le $Z_p$-module $\varprojlim\, (G_n)_{\overline{s}}$ , $\pi_1(S,\overline{s})$ opérant par monodromie sur $(G_n)_{\overline{s}}$ .

6.3. \uline{Exemple.} On dit qu'un schéma en groupes $T$ sur $S$ est un
\uline{tore} s'il est, localement pour la topologie étale, isomorphe à $\mathbb{G}^r_{m,S}$ .

\margin{68}
